% Written by moonshotai/kimi-k3.1-pro (Open Math OS), 2026-11-13; unreviewed AI draft.
\documentclass[10pt]{article}
\usepackage[margin=.85in]{geometry}
\usepackage{amsmath,amssymb,amsthm}
\usepackage{xcolor}
\usepackage[colorlinks=true,linkcolor=black]{hyperref}
\newtheorem{lemma}{Lemma}
\newtheorem{proposition}{Proposition}
\newtheorem{theorem}{Theorem}
\theoremstyle{definition}\newtheorem*{definition}{Definition}
\theoremstyle{remark}\newtheorem*{remark}{Remark}
\newcommand{\F}{\mathfrak F}
\newcommand{\GAP}[1]{\textcolor{red}{\textbf{GAP: }#1}}
\title{Every internal node is worth $-p$:\\the accounting proof for Fuss--Catalan numbers}
\author{Mighty Math Magicians --- unreviewed AI draft}
\date{2026-11-13}
\begin{document}\maketitle
\paragraph{Status.} This is an AI-written proof draft, not a human-signed theorem.
It derives a factorial valuation by symmetry; see Remark~\ref{qualification}.
Verification is in \texttt{experiments/E004\_accounting\_proof/}, with
qualifications in \texttt{GAPS\_accounting.md}. No \texttt{\char92 GAP} markers
are used.

Fix a prime $p$, $k\ge1$, and $q=p^k$. A $p$-ary tree is a leaf or an ordered
$p$-tuple of such trees. It has $N(m)=(m-1)/(p-1)$ internal nodes if it has $m$
leaves. Write $m=\sum_jm_jp^j$ in base $p$, $s=s_p(m)=\sum_jm_j$, and
\[
 \nu(m)=\frac{m-s}{p-1}=\sum_{j\ge1}\left\lfloor\frac m{p^j}\right\rfloor,
 \qquad t=N(m)-\nu(m)=\frac{s-1}{p-1}.
\]
The Fuss--Catalan number $C_n^{(q)}=\binom{qn}{n}/((q-1)n+1)$ counts ordered
full $q$-ary trees with $n$ nodes. Set $m=(q-1)n+1$. Expanding every $q$-ary
node into a complete $p$-ary block of height $k$ converts every such tree to
a $p$-ary tree with every leaf depth divisible by $k$.

\section{L1 and L2: labelled trees and cherries}
\begin{definition} In a \emph{leaf-labelled cyclic $k$-blocked $p$-ary tree},
every leaf has one of distinct labels $1,\ldots,m$ and lies at depth divisible
by $k$; at every internal node the $p$ children form a cyclic sequence. Its
\emph{rotation class} is the equivalence class of ordered trees obtained by
rotating any nonempty family of the child cycles.\end{definition}
\begin{lemma}[L1]\label{L1} A labelled rotation class has exactly $p^{N(m)}$
ordered representatives. If $L(m)$ is the number of classes, then
$m!C_n^{(q)}=p^{N(m)}L(m)$.\end{lemma}
\begin{proof} A local rotation acting trivially induces an automorphism fixing
all leaf labels. Recursively it fixes every node and child position, so it is
the identity. Hence all local rotations act freely on ordered labelled trees.
This exhibits every orbit as a product of one $\mathbb Z_p$ for each of $N(m)$
nodes. There are $m!C_n^{(q)}$ ordered labelled trees in total.\end{proof}

For general cyclic trees, each leaf also has \emph{class} $c\in\mathbb Z/k$;
admissibility requires its depth to be $\equiv-c\pmod k$.
\begin{lemma}\label{cherry} Cyclicly permuting $p$ distinct same-class labels
fixes precisely those admissible cyclic trees in which those labels form
$p$ sibling leaves. Collapsing them on both sides gives a bijection with $p-1$
copies of admissible trees with one fewer $p$-ary node and one class-$(c+1)$ leaf.\end{lemma}
\begin{proof} Let $\varphi$ be the induced automorphism and $v$ the lowest common
ancestor of the selected leaves. If $\varphi$ acts trivially at $v$, distinct
subtrees containing those leaves must be preserved; that is impossible because
their labels are cycled. A nonzero rotation then moves every leaf below $v$.
All other leaves are fixed, so the children of $v$ are precisely the selected
leaves. Once $x_1$ is placed in slot $1$, the other labels are forced by powers
$\sigma^u$ where $u=1,\ldots,p-1$. All $p-1$ choices are fixed. Conversely,
inverting collapse simply expands a class-$(c+1)$ leaf.\end{proof}

\begin{lemma}[L2]\label{L2} Let $\mathcal M(r)$ be the set of all unrestricted
cyclic p-ary trees with $r$ labelled leaves. Then
$|\mathcal M(r)|\equiv(-1)^{N(r)}\pmod p$, and hence
$L(m)\equiv(-1)^{N(m)}\pmod p$.\end{lemma}
\begin{proof} If $r=1$ the set is a singleton; otherwise $r\equiv1\pmod{p-1}$
is necessary. Apply Lemma~\ref{cherry} to fixed labels $1,\ldots,p$. Because
all orbits other than the fixed points have size $p$,
$|\mathcal M(r)|\equiv(p-1)|\mathcal M(r-p+1)|\pmod p$. Iterate exactly
$N(r)$ times. For blocked trees choose the next $p$ of the still present labels;
all class shifts remain zero at every collapse, giving the same recursion.\end{proof}

\section{L3: complete-tree forests}
\begin{definition} A \emph{digit-forest structure} on $[m]$ is a collection of
$m_j$ cyclic complete $p$-ary trees of height $j$, all disjoint, with leaf sets
partitioning $[m]$. Let $\F(m)$ denote the set of these structures. A tree of
height zero is a labelled leaf; the empty forest on $[0]$ is one object.\end{definition}
\begin{lemma}\label{stabilizer} $S_m$ is transitive on $\F(m)$, and its
stabilizer is $\prod_j(C_p^{N(p^j)}\wr S_{m_j})$, of order
$p^{\nu(m)}\prod_jm_j!$. Therefore
\begin{equation}\label{digit}
 m!=p^{\nu(m)}\prod_jm_j!\,|\F(m)|.
\end{equation}\end{lemma}
\begin{proof} Replacing labels one by one moves any listed leaf assignment to
another, and replacing entire equally high trees acts by separate root
permutations. Every tree's stabilizers come only from rotations at its $N(p^j)$
nodes; there cannot be an automorphism exchanging different subtrees and
preserving all labelled leaves. Since $\nu(m)=\sum m_jN(p^j)$, the order follows.\end{proof}

\begin{lemma}\label{K} Let $K(a)$ count permutations of $[pa]$ all of whose
cycles have length $p$. For $a\ge1$,
\begin{equation}\label{Kexact}
 K(a)=\binom{pa-1}{p-1}(p-1)!\,K(a-1)=\frac{(pa)!}{p^aa!},\qquad
 K(a)\equiv(-1)^a\pmod p.
\end{equation}\end{lemma}
\begin{proof} To count, choose the other members and the cyclic order of the
cycle containing $1$, then finish the remainder; this proves the exact
recursion and product formula. Conjugate the set by the cycle $(1\cdots p)$.
A fixed permutation preserves that cycle in one of $p-1$ ways and has arbitrary
$p$-cycle structure on the remaining points. The nonfixed orbits have size $p$,
so $K(a)\equiv(p-1)K(a-1)\equiv-K(a-1)$. With $K(0)=1$, the result follows.\end{proof}

\begin{lemma}\label{recursion} If $a=\lfloor m/p\rfloor$ and $r=m-pa$,
\begin{equation}\label{forest}
 |\F(m)|=\binom mrK(a)|\F(a)|,\qquad
 |\F(m)|\equiv(-1)^{\nu(m)}\pmod p.
\end{equation}\end{lemma}
\begin{proof} Exactly $r$ trees have height zero, so choose their leaves in
$\binom mr$ ways. Retain their $p$ cyclic complete subtrees of height $j$. These lower cycles
form one $p$-cycle-only permutation on the original labelled leaves, and the
retained upper forest is a digit forest with \(m_{j+1}\) trees of height
$j+1$ that became height $j$. Temporarily think of every removed root as an
additional leaf attached to the \(p\) existing child roots: that leaf has the
same tree's root sequence and can store its $p$-child group.

To label the new upper forest canonically, list its components increasingly
by their smallest original leaf and replace the retained components by their
roots. This assigns exactly the labels $1,\ldots,a$ to the retained roots.
The leaf lists remain cyclic and fixed under refitting the stored lower
cycles. The result has base-$p$ digits $a_j=m_{j+1}$ as required. The inverse
unlabels those \(a\) roots according to the lowest original leaf in each
$p$-cycle, and restores the \(p\)-tuple of complete subtrees at each stored
root. Because the collection of groups was disjoint and the child's cycles
were never altered, this is reconstructive. In particular the object being
counted $K(a)$ includes cyclic order and not just the partition into child
roots; no rotated equivalent is double counted. This completes \eqref{forest}.

For the congruence, rotate separately within each of $a$ consecutive blocks
of $p$ points of $[pa+r]$, holding the last $r$ points. A fixed $r$-subset
contains only entire moving blocks and fixed points. As $r<p$, there are no
complete blocks, so the fixed points form the unique fixed subset. Other
subset orbits of the $\mathbb Z_p$ action have size $p$. Hence
$\binom mr\equiv1\pmod p$; Lemma~\ref{K} and $\nu(m)=a+\nu(a)$ complete induction.\end{proof}

\section{The theorem}
\begin{theorem}\label{main} For every prime $p$, $k\ge1$ and $n\ge1$,
with $m=(p^k-1)n+1$,
\begin{equation}\label{maineq}
 \boxed{C_n^{(q)}\prod_jm_j!\equiv(-p)^t\pmod{p^{t+1}}}
\end{equation}
in $\mathbb Z$, and $v_p(C_n^{(q)})=t=(s_p(m)-1)/(p-1)$.\end{theorem}
\begin{proof} By Lemma~\ref{L1} and \eqref{digit},
\[
 p^{N(m)}L(m)=p^{\nu(m)}\prod_jm_j!|\F(m)|C_n^{(q)}.
\]
After cancelling common powers,
\[
 C_n^{(q)}\prod_jm_j!=p^t\frac{L(m)}{|\F(m)|}.
\]
The numerator and denominator are integral, well-defined, and $p$-units by
Lemmas~\ref{L2} and~\ref{recursion}. Hence $v_p(C_n^{(q)})=t$ because
each digit $m_j<p$ makes $\prod m_j!$ a $p$-unit. Modulo $p$ the fraction is
$(-1)^N/(-1)^\nu=(-1)^t$, proving \eqref{maineq}.\end{proof}
\begin{remark}\label{qualification} Here $v_p(m!)$ \emph{is} computed along the
way but from complete-tree symmetry rather than from Legendre's formula. Thus
this note does not meet an absolute prohibition on factorial valuations; whether
the alternative proof satisfies the problem's intended interpretation is for
the humans to judge. In the original skeleton proof, these digit trees are the
atoms, and the remaining $t$ internal nodes are the carry skeleton. A carry
replaces $p$ atoms of one height by one atom of the next height, read modulo
$k$; label cycling and collapsing translate that carry into Lemma~\ref{cherry}.\end{remark}

\section{Sand on a ring and on a line}
Take $a_c\ge0$ grains at sites $c=0,1,\ldots$ with finite support, or $c\in
\mathbb Z/k$. A site with $a_c\ge p$ topples to $(a_c-p,a_{c+1}+1)$, modulo $k$
on the ring. At the one-site ring, grain $1$ returns to the same site.
\begin{proposition}\label{sand} Every legal sequence is finite, and all maximal
legal sequences from a state have the same final state and per-site toppling
counts. From $m$ grains at the origin:
\begin{enumerate}
\item On the half-line the terminal state is the base-$p$ digits of $m$;
$x_c=\lfloor m/p^{c+1}\rfloor$ grains topple at site $c$, totalling $\nu(m)$.
\item On the ring the conserved quantity is $w(a)=\sum a_cp^c\pmod{q-1}$. The
terminal state is that of the unique $r\in\{1,\ldots,q-1\}$ congruent to $m$
modulo $q-1$ (with its $k$ digits zero-extended). For $q=2$ the final grain is
always $1$. For $m=(q-1)n+1$, the terminal state is $e_0$ and topplings total
$N(m)$.
\item Starting from the line's terminal digits folded into $k$ sites, exactly $t$
ring topplings remain. Therefore $v_p(C_n^{(q)})$ is the difference of the
topplings on the ring and on the line, both starting with $m$ at site zero.
\item If $S(a)_c=a_{c+1}$ on the ring, then $\mathbb Z^k/\operatorname{im}(pI-S)
\cong\mathbb Z/(q-1)\mathbb Z$.
\end{enumerate}\end{proposition}

\begin{proof} Each firing lowers total grains by $p-1$, hence every sequence
is finite. For abelianness, induct on the initial grain sum. Let $v,w$ be two
maximal legal words with first letters $c,d$. If $c=d$, remove them. If not,
other sites' firings only preserve or add grains at $c$. If $c$ never occurs in
$w$, it remains legal at its end, contradicting maximality. Before $c$'s first
occurrence in $w$, all firings are by other sites, so $c$ stays legal and each
commutes with $c$; commuting both orders leaves the other site's firing legal
because a preceding firing from a neighbour only supplies more grain. Move $c$
to the front. Remove the common first firing and invoke induction. This proves
same counts and final state.

The toppling count $x$ obeys
\[
 a-a_{\mathrm{final}}=(pI-S)x,
\]
with an extra site carrying grains past the known finite origin on the half-line.
The half-line preserves the integer $\sum a_cp^c=m$ exactly. A firing from site
$c$ loses $p$ grains there and gains one at $c+1$, so $m=(p-1)\sum x_cp^{c+1}
+\sum a_cp^c$. Modulo $p$, this gives $a_0=m_0$; subtract $a_0$, divide by $p$,
and repeat to identify every final digit. This also identifies the coefficients
$x_c$, since $x_c$ occurs with coefficient $(p-1)p^{c+1}$, and the inequality
$x_c\ge\sum_{i\ge c+1}m_ip^{i-c-1}$ follows from looking at grain flow as
an ordinary carry. (This conclusion is $x_c=
\lfloor m^{(c)}/p\rfloor$, where $m^{(0)}=m$ and $m^{(c+1)}
=(m^{(c)}-m_c)/p$; $\sum m_ip^{i-c-1}$ is its stable configuration.)
Consequently $\sum x_c=\nu(m)$.

On the ring the same weight is conserved modulo $q-1$ because $p^k\equiv1$.
The terminal digits are in $\{0,\ldots,p-1\}$, so their integer weight lies
between $0$ and $q-1$. A terminal state reached from positive grains has
weight not $0$ modulo $q-1$, hence its integer weight cannot be $0$. Its first
nonzero digit index is shifted by $k$ after a wraparound carry, so the canonical
weight representative $r$ in $1,\ldots,q-1$ gives precisely the stated final
state apart from index reposition. From the origin this index stays $0$ by the
abelian result, and $r=q-1$ if $q-1\mid m$. Since sum drops by $p-1$ the total
is $(m-s_p(r))/(p-1)$. For $m=(q-1)n+1$, $r=1$ and this is $N(m)$. Input
folded line digits has total $s_p(m)$, hence remaining $(s_p(m)-1)/(p-1)=t$.

For (4), $\det(pI-S)=p^k-1$ by expanding the permutation's nonzero term.
Deleting row $1$ and column $k$ gives a lower-triangular determinant: even for
$k=2$ this is $[-1]$, and $\det=\pm1$. The $(k-1)$-st determinant divisor of
the Smith normal form is $1$. By total determinant $|A|$ the final factor is
$p^k-1$, so the integer cokernel is cyclic of this order.\end{proof}
\end{document}
